\chapter{Andrei Okounkov --- Fields Medal 2006}\label{ch2006andreiokounkov:andrei-okounkov-fields-medal-2006-a-comprehensive-report-on-his-work-and-impact}

\begin{center}\rule{0.5\linewidth}{0.5pt}\end{center}

\textbf{Prepared:} July 2026 \textbf{Subject:} Andrei Okounkov (b. July
26, 1969, Moscow, Russia) \textbf{Position:} Professor of Mathematics,
Columbia University (since 2010) \textbf{Doctoral Advisor:} Alexander
Kirillov (Moscow State University) \textbf{Fields Medal Citation:}
\emph{``For his contributions to bridging probability, representation
theory and algebraic geometry.''}

\begin{center}\rule{0.5\linewidth}{0.5pt}\end{center}

\section{Table of Contents}\label{ch2006andreiokounkov:table-of-contents}

\begin{enumerate}
\def\labelenumi{\arabic{enumi}.}
\tightlist
\item
 Introduction
\item
 Biographical Background and Education
\item
 Random Partitions and the Plancherel Measure
\item
 The Asymptotic Representation Theory of Symmetric Groups
\item
 The Okounkov--Vershik Approach to the Symmetric Group
\item
 Gromov--Witten Invariants and Donaldson--Thomas Theory
\item
 The MNOP Conjecture
\item
 Nekrasov's Formula and the Seiberg--Witten Prepotential
\item
 Hilbert Schemes of Points on Surfaces
\item
 Quantum Integrable Systems
\item
 Okounkov Bodies in Algebraic Geometry
\item
 The Okounkov--Pandharipande Correspondence
\item
 Enumerative Geometry and Moduli Spaces
\item
 Tail Algebras and the Frozen Boundary
\item
 Representation Theory of Infinite-Dimensional Groups
\item
 Impact on Mathematics
\item
 Awards and Recognition
\item
 Selected Major Publications
\item
 Ongoing Work and Future Directions
\item
 Conclusion
\item
 References and Further Reading
\end{enumerate}

\begin{prereqs} Representation theory (symmetric groups, infinite-dimensional algebras), basic probability theory, complex analysis. \end{prereqs}

\begin{objectives} Explain the connection between random matrices and number theory; describe the representation theory of symmetric groups; understand the Borodin-Okounkov formula. \end{objectives}


\begin{center}\rule{0.5\linewidth}{0.5pt}\end{center}

\section{1. Introduction}\label{ch2006andreiokounkov:introduction}

Andrei Okounkov was awarded the Fields Medal in 2006 at the age of 36
for his notable work spanning probability, representation theory,
and algebraic geometry. The citation --- ``for his contributions to
bridging probability, representation theory and algebraic geometry'' ---
captures the defining character of his mathematics: the ability to
uncover deep connections between fields that had been thought separate,
and to use ideas from one to solve problems in another.

Okounkov's career has followed a distinctive arc. His early work on the
asymptotic representation theory of the symmetric groups, conducted
jointly with Grigori Vershik, transformed a classical subject by
introducing probabilistic and combinatorial methods. His later work on
enumerative geometry --- including the proof of the MNOP conjecture with
Rahul Pandharipande and his contributions to the Gromov--Witten /
Donaldson--Thomas correspondence --- placed him at the center of some of
the important developments in algebraic geometry of the early
twenty-first century.

This report provides a detailed account of Okounkov's major
contributions, the mathematical context in which they arose, and their
lasting impact.

\begin{center}\rule{0.5\linewidth}{0.5pt}\end{center}

\section{2. Biographical Background and
Education}\label{ch2006andreiokounkov:biographical-background-and-education}

\subsection{2.1 Early Life}\label{ch2006andreiokounkov:early-life}

Andrei Okounkov was born in Moscow in 1969. He demonstrated exceptional
mathematical ability from an early age and was selected for the
prestigious Kolmogorov boarding school (Specialized School No.~18)
affiliated with Moscow State University.

\subsection{2.2 University Education}\label{ch2006andreiokounkov:university-education}

Okounkov entered the Department of Mechanics and Mathematics at Moscow
State University in 1986. During his undergraduate years, he began
research under the supervision of Alexander Kirillov, working on
representation theory. He completed his PhD in 1995, writing a thesis on
``Asymptotic representation theory of the symmetric group.''

\subsection{2.3 Academic Career}\label{ch2006andreiokounkov:academic-career}

After completing his graduate studies, Okounkov held a series of
prestigious positions: - \textbf{1995--1998:} Postdoctoral positions at
the Institute for Advanced Study (Princeton), University of Chicago, and
IHP (Paris). - \textbf{1998--2002:} Assistant/Associate Professor,
University of Chicago. - \textbf{2002--2009:} Professor, Princeton
University. - \textbf{2009--2010:} Professor, University of California,
Berkeley. - \textbf{2010--present:} Professor, Columbia University.

He has also held visiting professorships at the Institute for Advanced
Study, the Kavli Institute for Theoretical Physics, and IHES.

\begin{center}\rule{0.5\linewidth}{0.5pt}\end{center}

\section{3. Random Partitions and the Plancherel
Measure}\label{ch2006andreiokounkov:random-partitions-and-the-plancherel-measure}

\subsection{3.1 Partitions and the Symmetric
Group}\label{ch2006andreiokounkov:partitions-and-the-symmetric-group}

A partition
\(\lambda = (\lambda_1 \ge \lambda_2 \ge \cdots \ge \lambda_\ell > 0)\)
is a weakly decreasing sequence of positive integers summing to \(n\).
Partitions index both the conjugacy classes and the irreducible
representations of the symmetric group \(S_n\). The \textbf{Plancherel
measure} on partitions of \(n\) assigns to a partition \(\lambda\) the
probability:

\[
\mathbb{P}(\lambda) = \frac{(\dim \lambda)^2}{n!}
\]

where \(\dim \lambda\) is the dimension of the corresponding irreducible
representation. This measure arises naturally from the Fourier transform
on \(S_n\) and from the decomposition of the regular representation.

\subsection{3.2 The Limit Shape of Random
Partitions}\label{ch2006andreiokounkov:the-limit-shape-of-random-partitions}

The study of the Plancherel measure for large \(n\) is a central problem
in asymptotic representation theory. A notable discovery, made
independently by Okounkov, by Richard Stanley, and by Boris Kerov and
Anatoly Vershik, is that a random partition under the Plancherel measure
exhibits a \textbf{limit shape}.

If we view a partition as a Young diagram (with rows of length
\(\lambda_i\)), and rescale by a factor of \(1/\sqrt{n}\), the boundary
of the diagram converges almost surely to a deterministic curve:

\[
\Omega(x) = \begin{cases}
\frac{2}{\pi} \left( x \arcsin\frac{x}{2} + \sqrt{4 - x^2} \right) & |x| \le 2 \\
|x| & |x| > 2
\end{cases}
\]

This is the \textbf{Vershik--Kerov--Logan--Shepp} limit shape, and it
coincides with the limiting distribution of eigenvalues of random
matrices from the Gaussian Unitary Ensemble (GUE).

\subsection{3.3 Okounkov's Contributions}\label{ch2006andreiokounkov:okounkovs-contributions}

Okounkov made fundamental contributions to the study of the Plancherel
measure, including:

\begin{enumerate}
\def\labelenumi{\arabic{enumi}.}
\tightlist
\item
 \textbf{Explicit asymptotics:} Exact formulas for the correlation
 functions of the Plancherel measure in terms of the discrete sine
 kernel.
\item
 \textbf{Relation to random matrices:} A precise correspondence between
 the fluctuations of random partitions and the eigenvalues of random
 matrices.
\item
 \textbf{The Arcsine law:} The distribution of the length of the first
 row of a random partition follows the Tracy--Widom distribution from
 random matrix theory.
\end{enumerate}

These results established the Plancherel measure as a fundamental model
of random combinatorial structures, on par with random matrices and
random permutations.

\begin{center}\rule{0.5\linewidth}{0.5pt}\end{center}

\section{4. The Asymptotic Representation Theory of Symmetric
Groups}\label{ch2006andreiokounkov:the-asymptotic-representation-theory-of-symmetric-groups}

\subsection{4.1 The Problem}\label{ch2006andreiokounkov:the-problem}

The representation theory of the symmetric group \(S_n\) for fixed \(n\)
is classical (Frobenius, Young, Schur). The problem of understanding the
behavior of representations as \(n \to \infty\) --- \textbf{asymptotic
representation theory} --- was initiated by Vershik and Kerov in the
1970s.

The central question: consider a sequence of irreducible representations
of \(S_n\) whose dimensions grow appropriately with \(n\). What are the
statistical properties of such representations? How do the characters
behave in the limit?

\subsection{4.2 Okounkov's
Contributions}\label{ch2006andreiokounkov:okounkovs-contributions-1}

Okounkov, building on his work with Vershik, established the foundations
of asymptotic representation theory:

\begin{enumerate}
\def\labelenumi{\arabic{enumi}.}
\tightlist
\item
 \textbf{The character formula:} New formulas for the characters of the
 symmetric group in terms of the geometry of Young diagrams.
\item
 \textbf{The Central Limit Theorem:} Fluctuations of characters around
 their limits are Gaussian.
\item
 \textbf{Free probability:} Connections between the asymptotic
 representation theory of \(S_n\) and Voiculescu's free probability
 theory.
\end{enumerate}

This work revealed deep connections between representation theory and
probability theory, showing that the large-\(n\) limit of the symmetric
group exhibits striking parallels with random matrix theory and free
probability.

\begin{center}\rule{0.5\linewidth}{0.5pt}\end{center}

\section{5. The Okounkov--Vershik Approach to the Symmetric
Group}\label{ch2006andreiokounkov:the-okounkovvershik-approach-to-the-symmetric-group}

\subsection{5.1 A New Approach}\label{ch2006andreiokounkov:a-new-approach}

Traditionally, the irreducible representations of \(S_n\) are
constructed using Young symmetrizers and the Young--Jucys--Murphy
elements. In a landmark 1996 paper, Okounkov and Vershik proposed a
radically different approach based on:

\begin{enumerate}
\def\labelenumi{\arabic{enumi}.}
\item
 \textbf{The Gelfand--Tsetlin basis:} Starting from the chain of
 subgroups: \[
 S_1 \subset S_2 \subset \cdots \subset S_n
 \] they construct representations by analyzing the branching rules at
 each step.
\item
 \textbf{The Young graph:} The set of all partitions, with edges
 connecting \(\lambda \vdash n\) to \(\mu \vdash (n+1)\) when \(\mu\)
 is obtained from \(\lambda\) by adding a box. This is the
 \textbf{Young graph} or the \textbf{Pascal graph of partitions}.
\item
 \textbf{Central measures:} Classification of all possible central
 measures on the Young graph --- that is, measures that are consistent
 under the branching rules.
\end{enumerate}

\subsection{5.2 The Main Theorem}\label{ch2006andreiokounkov:the-main-theorem}

Okounkov and Vershik proved that the set of all \emph{central measures}
on the Young graph is in bijection with a certain set of parameters ---
the \textbf{Thoma simplex} --- which had earlier appeared in the work of
Thoma on characters of the infinite symmetric group.

They also gave a new construction of the irreducible representations of
\(S_n\) using the Young--Jucys--Murphy elements, and they proved that
the Bratteli diagram of the symmetric groups is isomorphic to the Young
graph.

This work provided the definitive modern framework for the
representation theory of the symmetric groups and their
infinite-dimensional limits.

\subsection{5.3 Impact}\label{ch2006andreiokounkov:impact}

The Okounkov--Vershik approach is now the standard way to study the
representation theory of symmetric groups and their affine
generalizations. It has been extended to: - \textbf{Hecke algebras:} The
\(q\)-deformed version of the symmetric group. - \textbf{Affine Lie
algebras:} The representation theory of \(\widehat{\mathfrak{sl}}_n\). -
\textbf{Double affine Hecke algebras:} The representation theory of
Cherednik's DAHA. - \textbf{Wreath products:} Symmetric groups with
additional structure.

\begin{center}\rule{0.5\linewidth}{0.5pt}\end{center}

\section{6. Gromov--Witten Invariants and Donaldson--Thomas
Theory}\label{ch2006andreiokounkov:gromovwitten-invariants-and-donaldsonthomas-theory}

\subsection{6.1 Enumerative Geometry}\label{ch2006andreiokounkov:enumerative-geometry}

\textbf{Enumerative geometry} counts geometric objects satisfying
incidence conditions. The classical problem --- ``How many rational
curves of degree \(d\) pass through \(3d-1\) points in the plane?'' ---
was solved in the nineteenth century but resisted rigorous treatment
until the advent of \textbf{Gromov--Witten theory}.

Gromov--Witten invariants count (pseudo-)holomorphic curves in a smooth
projective variety, and they are defined via intersection theory on the
moduli space of stable maps. Donaldson--Thomas invariants, by contrast,
count sheaves with fixed numeric invariants, and they are defined via
the moduli space of ideal sheaves of curves.

\subsection{6.2 The MNOP Conjecture}\label{ch2006andreiokounkov:the-mnop-conjecture}

In 2003, Davesh Maulik, Nikita Nekrasov, Andrei Okounkov, and Rahul
Pandharipande (hence ``MNOP'') proposed a deep conjecture relating
Gromov--Witten invariants and Donaldson--Thomas invariants.

The conjecture states that for a nonsingular projective threefold \(X\),
the generating functions of Gromov--Witten invariants and
Donaldson--Thomas invariants are related by a simple change of
variables:

\[
Z_{GW}(X; q) = Z_{DT}(X; -q)^{\chi(X)}
\]

where \(\chi(X)\) is the topological Euler characteristic.

This relationship was completely unexpected. Gromov--Witten invariants
are defined in terms of curves (1-dimensional objects), while
Donaldson--Thomas invariants are defined in terms of sheaves
(0-dimensional objects). The MNOP conjecture asserted a deep duality
between these two worlds.

\subsection{6.3 Proof of the MNOP
Conjecture}\label{ch2006andreiokounkov:proof-of-the-mnop-conjecture}

Over the next several years, Okounkov and Pandharipande proved the MNOP
conjecture for a wide class of threefolds: - \textbf{Local curves:} The
total space of a vector bundle over a curve. - \textbf{Toric
threefolds:} Using the localization techniques of equivariant geometry.
- \textbf{Calabi--Yau threefolds:} Under certain conditions.

The proof uses the method of \textbf{virtual localization} on the moduli
space of stable pairs (Pandharipande--Thomas theory), combined with
sophisticated combinatorial analysis of the fixed point contributions.

\begin{center}\rule{0.5\linewidth}{0.5pt}\end{center}

\section{7. The Gromov--Witten / Donaldson--Thomas
Correspondence}\label{ch2006andreiokounkov:the-gromovwitten-donaldsonthomas-correspondence}

\subsection{7.1 The Correspondence}\label{ch2006andreiokounkov:the-correspondence}

The MNOP conjecture implies a far-reaching correspondence between two
different approaches to enumerative geometry:

{\def\LTcaptype{none} % do not increment counter
\begin{longtable}[]{@{}ll@{}}
\toprule\noalign{}
Gromov--Witten Theory & Donaldson--Thomas Theory \\
\midrule\noalign{}
\endhead
\bottomrule\noalign{}
\endlastfoot
Curves & Sheaves \\
Moduli of stable maps & Moduli of ideal sheaves \\
Virtual fundamental class & Virtual fundamental class \\
Descendent invariants & Descendent invariants \\
Connected invariants & Disconnected invariants \\
Genus expansion & No genus expansion \\
\end{longtable}
}

The correspondence has been verified for many important examples,
including all toric threefolds, local curves, and certain complete
intersections.

\subsection{7.2 Okounkov's
Contributions}\label{ch2006andreiokounkov:okounkovs-contributions-2}

Okounkov's work on the correspondence includes:

\begin{enumerate}
\def\labelenumi{\arabic{enumi}.}
\item
 \textbf{Proof for toric threefolds:} With Pandharipande, they proved
 the full Gromov--Witten / Donaldson--Thomas correspondence for toric
 threefolds by computing both sides using localization.
\item
 \textbf{The descendent correspondence:} Extension of the
 correspondence to include descendent invariants (insertions of Chern
 classes).
\item
 \textbf{Pandharipande--Thomas theory:} Work on the moduli space of
 stable pairs, which provides an alternative approach to
 Donaldson--Thomas theory.
\end{enumerate}

\subsection{7.3 Impact}\label{ch2006andreiokounkov:impact-1}

The MNOP conjecture and its proof transformed enumerative geometry. It
established Donaldson--Thomas theory as a fundamental tool in algebraic
geometry and revealed a deep unity between seemingly different
approaches to curve counting. The techniques developed --- particularly
the use of localization and the analysis of moduli spaces --- have been
applied to many other problems.

\begin{center}\rule{0.5\linewidth}{0.5pt}\end{center}

\section{8. Nekrasov's Formula and the Seiberg--Witten
Prepotential}\label{ch2006andreiokounkov:nekrasovs-formula-and-the-seibergwitten-prepotential}

\subsection{8.1 Background}\label{ch2006andreiokounkov:background}

In theoretical physics, the \textbf{Seiberg--Witten prepotential} is the
effective action of \(N = 2\) supersymmetric gauge theories in four
dimensions. It encodes the low-energy physics of the theory, including
the masses of BPS states and the geometry of the moduli space of vacua.

In 2002, Nikita Nekrasov proposed a formula for the instanton
contributions to the Seiberg--Witten prepotential of \(N = 2\)
supersymmetric Yang--Mills theory. The formula expresses the partition
function as:

\[
Z(\epsilon_1, \epsilon_2, a) = \sum_{k=0}^\infty \Lambda^{2k} \int_{M_{k,N}} 1
\]

where \(M_{k,N}\) is the moduli space of instantons and the integral is
over the virtual fundamental class.

\subsection{8.2 Nekrasov's Conjecture and Okounkov's
Proof}\label{ch2006andreiokounkov:nekrasovs-conjecture-and-okounkovs-proof}

Nekrasov conjectured that in the limit \(\epsilon_1, \epsilon_2 \to 0\),
his partition function reproduces the Seiberg--Witten prepotential:

\[
\lim_{\epsilon_1, \epsilon_2 \to 0} \epsilon_1 \epsilon_2 \log Z(\epsilon_1, \epsilon_2, a) = \mathcal{F}_{SW}(a)
\]

In 2003, Nekrasov and Okounkov published a complete proof of this
conjecture for \(U(N)\) gauge groups. The proof used:

\begin{enumerate}
\def\labelenumi{\arabic{enumi}.}
\item
 \textbf{The ADHM construction:} A description of the instanton moduli
 space as a hyper-Kähler quotient, which allows its virtual fundamental
 class to be computed.
\item
 \textbf{Equivariant localization:} The use of torus actions to reduce
 the integrals to sums over fixed points, indexed by pairs of
 partitions.
\item
 \textbf{The Cauchy identity:} A combinatorial identity for sums over
 partitions, which simplifies the resulting expression to the
 Seiberg--Witten form.
\end{enumerate}

The proof was a masterpiece of algebraic geometry and representation
theory, showing that the instanton partition function can be expressed
in terms of the representation theory of \(\mathfrak{gl}(\infty)\).

\subsection{8.3 Physicists' Formula}\label{ch2006andreiokounkov:physicists-formula}

The Nekrasov--Okounkov proof gave a rigorous foundation to a central
result in supersymmetric gauge theory. The formula --- known as the
\textbf{Nekrasov formula} or the \textbf{Nekrasov--Okounkov formula} ---
is now a standard tool in both physics and mathematics, with
applications to:

\begin{itemize}
\tightlist
\item
 \textbf{AGT correspondence:} The relationship between 4D gauge
 theories and 2D conformal field theories.
\item
 \textbf{Wall-crossing:} The transformation of BPS invariants under
 changes of stability conditions.
\item
 \textbf{Topological strings:} The relationship between gauge theory
 and string theory.
\end{itemize}

\begin{center}\rule{0.5\linewidth}{0.5pt}\end{center}

\section{9. Hilbert Schemes of Points on
Surfaces}\label{ch2006andreiokounkov:hilbert-schemes-of-points-on-surfaces}

\subsection{9.1 Definition}\label{ch2006andreiokounkov:definition}

The \textbf{Hilbert scheme of \(n\) points} on a smooth algebraic
surface \(S\), denoted \(S^{[n]}\), parametrizes zero-dimensional
subschemes of \(S\) of length \(n\). It is a resolution of the symmetric
product \(S^{(n)} = S^n / S_n\).

Hilbert schemes of points are fundamental objects in algebraic geometry,
with connections to: - \textbf{Moduli spaces:} They are the simplest
moduli spaces of sheaves. - \textbf{Representation theory:} They carry
actions of Heisenberg algebras. - \textbf{Topology:} Their cohomology is
known to be infinite-dimensional and to carry a rich algebraic
structure.

\subsection{9.2 Okounkov's
Contributions}\label{ch2006andreiokounkov:okounkovs-contributions-3}

Okounkov made several fundamental contributions to the theory of Hilbert
schemes:

\begin{enumerate}
\def\labelenumi{\arabic{enumi}.}
\item
 \textbf{Quantum cohomology:} With Davesh Maulik, he computed the
 quantum cohomology ring of \(S^{[n]}\) in terms of the quantum
 cohomology of \(S\).
\item
 \textbf{Stable envelopes:} The development of the theory of stable
 envelopes, which provide a systematic way to relate the equivariant
 cohomology of different moduli spaces.
\item
 \textbf{Elliptic stable envelopes:} A generalization of stable
 envelopes to the elliptic setting, which has applications to the
 theory of elliptic quantum groups.
\item
 \textbf{Hilbert schemes and K-theory:} New results on the K-theory of
 Hilbert schemes and their connections to vertex algebras.
\end{enumerate}

\subsection{9.3 The Maulik--Okounkov
Theory}\label{ch2006andreiokounkov:the-maulikokounkov-theory}

In a series of papers with Davesh Maulik, Okounkov developed a general
theory of \textbf{quantum cohomology} for moduli spaces. The key insight
is that the quantum cohomology of a moduli space can be computed using
the geometry of correspondences and the combinatorics of stable
envelopes.

The Maulik--Okounkov theory has been influential in: -
\textbf{Enumerative geometry:} A systematic approach to computing
Gromov--Witten invariants. - \textbf{Geometric representation theory:}
Connections to Yangians, quantum groups, and affine algebras. -
\textbf{Integrable systems:} The relationship between quantum cohomology
and quantum integrable models.

\begin{center}\rule{0.5\linewidth}{0.5pt}\end{center}

\section{10. Quantum Integrable
Systems}\label{ch2006andreiokounkov:quantum-integrable-systems}

\subsection{10.1 Integrability}\label{ch2006andreiokounkov:integrability}

A \textbf{quantum integrable system} is a quantum mechanical system with
as many commuting conserved quantities as degrees of freedom. The most
famous examples are: - The quantum Heisenberg spin chain. - The XXZ and
XYZ spin chains. - The Gaudin model.

\subsection{10.2 Okounkov's Contributions to Integrable
Systems}\label{ch2006andreiokounkov:okounkovs-contributions-to-integrable-systems}

Okounkov has made important contributions to the theory of quantum
integrable systems, including:

\begin{enumerate}
\def\labelenumi{\arabic{enumi}.}
\item
 \textbf{The quantum Toda lattice:} A complete solution of the quantum
 Toda lattice using representation theory of affine Lie algebras.
\item
 \textbf{Quantum cohomology and Bethe ansatz:} A connection between the
 quantum cohomology of the Grassmannian and the Bethe ansatz for the
 XXZ spin chain.
\item
 \textbf{Applied to Hilbert schemes:} The quantum cohomology of Hilbert
 schemes of points on surfaces is related to an integrable system ---
 the Calogero--Sutherland model.
\item
 \textbf{Stable envelopes and R-matrices:} A geometric construction of
 the R-matrix for quantum groups using stable envelopes on moduli
 spaces.
\end{enumerate}

\subsection{10.3 The Geometric R-Matrix}\label{ch2006andreiokounkov:the-geometric-r-matrix}

One of Okounkov's influential recent ideas is the \textbf{geometric
construction of the R-matrix} using stable envelopes. In joint work with
Maulik and others, he showed that the stable envelopes of a moduli space
satisfy a version of the Yang--Baxter equation, and thus give rise to an
R-matrix for a quantum group.

This construction has been applied to produce R-matrices for: - The
Yangian of \(\mathfrak{gl}_n\). - The quantum loop algebra of
\(\mathfrak{gl}_n\). - The elliptic quantum group of
\(\mathfrak{gl}_n\).

The geometric R-matrix provides a bridge between enumerative geometry
and representation theory, allowing geometric methods to be used to
study the representation theory of quantum groups.

\begin{center}\rule{0.5\linewidth}{0.5pt}\end{center}

\section{11. Okounkov Bodies in Algebraic
Geometry}\label{ch2006andreiokounkov:okounkov-bodies-in-algebraic-geometry}

\subsection{11.1 The Problem}\label{ch2006andreiokounkov:the-problem-1}

A fundamental problem in algebraic geometry is to understand the
asymptotic behavior of linear systems on a projective variety. Given a
line bundle \(L\) on a variety \(X\), the \textbf{volume} of \(L\) is:

\[
\operatorname{vol}(L) = \limsup_{m \to \infty} \frac{h^0(X, L^{\otimes m})}{m^n / n!}
\]

where \(n = \dim X\). The volume measures the ``size'' of the linear
system of \(L\), and it plays a central role in the theory of birational
geometry.

\subsection{11.2 Okounkov Bodies}\label{ch2006andreiokounkov:okounkov-bodies}

In 2003, Okounkov introduced a construction that associates to a line
bundle \(L\) a convex body \(\Delta(L) \subset \mathbb{R}^n\) --- now
called an \textbf{Okounkov body} --- whose Euclidean volume equals the
volume of \(L\).

The construction works by choosing a \textbf{flag} of subvarieties:

\[
X = X_0 \supset X_1 \supset \cdots \supset X_n = \{pt\}
\]

and then using the orders of vanishing of sections along the flag to map
the space of sections \(H^0(X, L^{\otimes m})\) to a semigroup in
\(\mathbb{Z}^n\). The closure of the limiting shape of this semigroup is
the Okounkov body.

\subsection{11.3 Applications}\label{ch2006andreiokounkov:applications}

Okounkov bodies have become a fundamental tool in birational geometry,
with applications including:

\begin{enumerate}
\def\labelenumi{\arabic{enumi}.}
\item
 \textbf{Theorem of Fujita:} A new proof of Fujita's approximation
 theorem for volumes.
\item
 \textbf{Concavity of volumes:} Khovanskii's theorem on the
 log-concavity of volumes.
\item
 \textbf{Zariski decompositions:} The existence of Zariski
 decompositions for line bundles on surfaces can be understood in terms
 of the Okounkov body.
\item
 \textbf{Newton--Okounkov bodies:} A generalization of the classical
 Newton polytope to the setting of algebraic varieties, developed
 jointly with Kaveh and Khovanskii.
\item
 \textbf{Seshadri constants:} Upper bounds on Seshadri constants in
 terms of the geometry of Okounkov bodies.
\end{enumerate}

Okounkov bodies have also been used to prove the
\textbf{Khovanskii--Teissier conjecture} about the concavity of
intersection numbers, and they have been applied to the study of mirror
symmetry and toric degenerations.

\begin{center}\rule{0.5\linewidth}{0.5pt}\end{center}

\section{12. The Okounkov--Pandharipande
Correspondence}\label{ch2006andreiokounkov:the-okounkovpandharipande-correspondence}

\subsection{12.1 The Idea}\label{ch2006andreiokounkov:the-idea}

The Okounkov--Pandharipande correspondence is a general principle
relating the enumeration of curves on a threefold to the representation
theory of infinite-dimensional Lie algebras. Specifically, the
Gromov--Witten invariants of a threefold \(X\) can be expressed as the
matrix elements of certain operators on a Fock space.

\subsection{12.2 The Fock Space Model}\label{ch2006andreiokounkov:the-fock-space-model}

A \textbf{Fock space} is a Hilbert space spanned by partitions, with
creation and annihilation operators that add or remove boxes from a
Young diagram. The fundamental observation of Okounkov and Pandharipande
is that the moduli space of curves on a threefold carries a natural
action of the Heisenberg algebra, with the creation operators
corresponding to curves of higher genus.

The generating function of Gromov--Witten invariants becomes:

\[
Z_{GW}(X; q) = \langle \text{vacuum} | \prod_i \exp(\text{operator}_i) | \text{vacuum} \rangle
\]

This is the \textbf{vacuum expectation value} of a product of operators
in the Fock space.

\subsection{12.3 The Virasoro
Constraints}\label{ch2006andreiokounkov:the-virasoro-constraints}

The Okounkov--Pandharipande correspondence also predicts that the
Gromov--Witten generating function satisfies a system of differential
equations --- the \textbf{Virasoro constraints} --- which are analogous
to the Ward identities in conformal field theory.

For a threefold, these constraints take the form:

\[
L_n Z_{GW} = 0, \quad n \ge -1
\]

where \(L_n\) are Virasoro operators acting on the Fock space. These
constraints completely determine the Gromov--Witten invariants once the
initial data is known.

\subsection{12.4 Impact}\label{ch2006andreiokounkov:impact-2}

The Okounkov--Pandharipande correspondence has had a substantial
effect on enumerative geometry. It has been used to: - Compute
Gromov--Witten invariants for many threefolds. - Prove the Virasoro
conjecture for target curves. - Connect Gromov--Witten theory to the
representation theory of infinite-dimensional algebras.

\begin{center}\rule{0.5\linewidth}{0.5pt}\end{center}

\section{13. Enumerative Geometry and Moduli
Spaces}\label{ch2006andreiokounkov:enumerative-geometry-and-moduli-spaces}

\subsection{13.1 The Moduli Space of
Curves}\label{ch2006andreiokounkov:the-moduli-space-of-curves}

The moduli space \(M_g\) of genus \(g\) curves is one of the most
important objects in algebraic geometry. Okounkov has contributed to the
understanding of its tautological ring and its intersection theory.

\subsection{13.2 Contributions}\label{ch2006andreiokounkov:contributions}

With Pandharipande and others, Okounkov made contributions to:

\begin{enumerate}
\def\labelenumi{\arabic{enumi}.}
\item
 \textbf{The ELSV formula:} A formula relating simple Hurwitz numbers
 to intersection numbers on \(M_g\), which was conjectured by Ekedahl,
 Lando, Shapiro, and Vainshtein, and proved by Okounkov and
 Pandharipande.
\item
 \textbf{The \(r\)-spin theory:} A generalization of Gromov--Witten
 theory to the setting of \(r\)-spin curves, with applications to the
 representation theory of \(W\)-algebras.
\item
 \textbf{Cohomological field theories:} The construction of
 cohomological field theories from Gromov--Witten invariants and their
 classification.
\end{enumerate}

\subsection{13.3 The ELSV Formula}\label{ch2006andreiokounkov:the-elsv-formula}

The ELSV formula states that:

\[
H_g(\mu) = \frac{(2g - 2 + n + |\mu|)!}{\prod_i \mu_i} \int_{\overline{M}_{g,n}} \frac{\Lambda_g^\vee(1)}{\prod_i (1 - \mu_i \psi_i)}
\]

where \(H_g(\mu)\) are simple Hurwitz numbers (counting branched covers
of \(\mathbb{P}^1\) with a given ramification profile) and
\(\Lambda_g^\vee\) is the Chern polynomial of the Hodge bundle. This
formula was conjectured by Ekedahl, Lando, Shapiro, and Vainshtein and
proved by Okounkov and Pandharipande using methods from Gromov--Witten
theory.

\begin{center}\rule{0.5\linewidth}{0.5pt}\end{center}

\section{14. Tail Algebras and the Frozen
Boundary}\label{ch2006andreiokounkov:tail-algebras-and-the-frozen-boundary}

\subsection{14.1 Random Surfaces}\label{ch2006andreiokounkov:random-surfaces}

A recurring theme in Okounkov's work is the study of \textbf{random
surfaces} --- random interfaces in statistical mechanics models such as
the dimer model, the six-vertex model, and the Ising model.

\subsection{14.2 The Frozen Boundary}\label{ch2006andreiokounkov:the-frozen-boundary}

In joint work with Richard Kenyon and Scott Sheffield, Okounkov studied
the \textbf{dimer model} on the hexagonal lattice, which is equivalent
to the problem of random tilings of a region by lozenges (rhombi).

They proved that for a large class of regions, the random tiling
exhibits a \textbf{frozen boundary} --- a curve that separates a region
of high entropy (where all three orientations of lozenges appear) from a
region of low entropy (where only two orientations appear). This frozen
boundary is described by an algebraic curve, and its shape depends on
the geometry of the region.

This work has deep connections to: - \textbf{Algebraic geometry:} The
frozen boundary is a real algebraic curve. - \textbf{Complex analysis:}
The shape of the boundary is determined by the Schwarz--Christoffel
mapping. - \textbf{Probability theory:} The fluctuations of the
interface are described by the Gaussian free field.

\subsection{14.3 The Limit Shape}\label{ch2006andreiokounkov:the-limit-shape}

Okounkov and his collaborators showed that, as the size of the region
tends to infinity, the random tiling converges almost surely to a
\textbf{limit shape} --- a surface that minimizes a certain surface
tension functional. The frozen boundary is the boundary of the ``liquid
region'' where the surface is curved.

This was a spectacular confirmation of the general theory of random
interfaces and a testament to the power of integrable probability.

\begin{center}\rule{0.5\linewidth}{0.5pt}\end{center}

\section{15. Representation Theory of Infinite-Dimensional
Groups}\label{ch2006andreiokounkov:representation-theory-of-infinite-dimensional-groups}

\subsection{15.1 The Infinite Symmetric
Group}\label{ch2006andreiokounkov:the-infinite-symmetric-group}

The \textbf{infinite symmetric group} \(S_\infty = \bigcup_n S_n\) is
the group of finitary permutations of a countable set. Its
representation theory is governed by the work of Thoma, Vershik, and
Kerov, and it is intimately connected to the asymptotic representation
theory of the finite symmetric groups.

\subsection{15.2 Okounkov's
Contributions}\label{ch2006andreiokounkov:okounkovs-contributions-4}

Okounkov studied the representation theory of \(S_\infty\) from the
perspective of the Okounkov--Vershik approach. He showed:

\begin{enumerate}
\def\labelenumi{\arabic{enumi}.}
\item
 \textbf{Classification of characters:} All characters of \(S_\infty\)
 are parametrized by Thoma parameters \((\alpha, \beta, \delta)\) where
 \(\alpha_i, \beta_i\) are nonnegative sequences summing to at most
 \(\delta\).
\item
 \textbf{The Plancherel measure:} The Plancherel measure on partitions
 of \(n\) converges to the \textbf{Vershik--Kerov measure} on Thoma
 parameters as \(n \to \infty\).
\item
 \textbf{Free probability:} The characters of \(S_\infty\) satisfy the
 same formulas as the cumulants in free probability theory.
\end{enumerate}

\begin{center}\rule{0.5\linewidth}{0.5pt}\end{center}

\section{16. Impact on Mathematics}\label{ch2006andreiokounkov:impact-on-mathematics}

\subsection{16.1 Impact on Enumerative
Geometry}\label{ch2006andreiokounkov:impact-on-enumerative-geometry}

Okounkov's work on Gromov--Witten invariants, Donaldson--Thomas theory,
and the MNOP conjecture transformed enumerative geometry. The MNOP
conjecture revealed a deep unity between two apparently different
counting theories, and the techniques developed through its proof ---
virtual localization, the Fock space approach, the Virasoro constraints
--- have become standard tools.

\subsection{16.2 Impact on Representation
Theory}\label{ch2006andreiokounkov:impact-on-representation-theory}

The Okounkov--Vershik approach to the symmetric group provided the
definitive modern framework for the subject. It clarified the role of
the Young--Jucys--Murphy elements, unified the different constructions
of the irreducible representations, and opened the door to the
asymptotic theory. The geometric construction of R-matrices using stable
envelopes is a major contribution to geometric representation theory.

\subsection{16.3 Impact on Probability}\label{ch2006andreiokounkov:impact-on-probability}

Okounkov's work on the Plancherel measure, random partitions, and random
surfaces established deep connections between probability theory and
algebraic combinatorics. The limit shape of random partitions, the
frozen boundary of random tilings, and the Tracy--Widom distribution of
fluctuations are now central objects in integrable probability.

\subsection{16.4 Impact on Algebraic
Geometry}\label{ch2006andreiokounkov:impact-on-algebraic-geometry}

Okounkov bodies have become a standard tool in birational geometry. They
provide a convex-geometric approach to volumes, positivity, and
asymptotic invariants of line bundles, and they have been used to prove
fundamental results about the geometry of algebraic varieties.

\subsection{16.5 Impact on Mathematical
Physics}\label{ch2006andreiokounkov:impact-on-mathematical-physics}

The Nekrasov--Okounkov proof of Nekrasov's conjecture provided rigorous
foundations for the Seiberg--Witten prepotential and its relation to
instanton counting. The AGT correspondence, which relates 4D gauge
theories to 2D conformal field theories, relies crucially on the
Nekrasov--Okounkov formula.

\begin{center}\rule{0.5\linewidth}{0.5pt}\end{center}

\section{17. Awards and Recognition}\label{ch2006andreiokounkov:awards-and-recognition}

\textbf{2006:} - Fields Medal (International Mathematical Union) -
European Prize in Combinatorics

\textbf{2004:} - Invited Speaker, International Congress of
Mathematicians (Beijing) - Clay Research Scholar (Clay Mathematics
Institute)

\textbf{2003:} - EMS Prize (European Mathematical Society)

\textbf{2000:} - Packard Fellowship for Science and Engineering

\textbf{1998--2000:} - Clay Research Fellow

\begin{center}\rule{0.5\linewidth}{0.5pt}\end{center}

\section{18. Selected Major
Publications}\label{ch2006andreiokounkov:selected-major-publications}

\begin{enumerate}
\def\labelenumi{\arabic{enumi}.}
\item
 \textbf{``Asymptotics of the Plancherel measure for the symmetric
 group''} (with S. Kerov, G. Olshanski, A. Vershik). \emph{Journal of
 Mathematical Sciences} 87(5), 3806--3856, 1997. --- \emph{The limit
 shape of random partitions and the distribution of their
 fluctuations.}
\item
 \textbf{``A new approach to the representation theory of the symmetric
 groups''} (with A. Vershik). \emph{Selecta Mathematica} 2(4),
 581--605, 1996. --- \emph{The Okounkov--Vershik approach to the
 symmetric group, providing a new construction of irreducible
 representations.}
\item
 \textbf{``The equivalence of Gromov--Witten and Donaldson--Thomas
 invariants''} (with D. Maulik, N. Nekrasov, R. Pandharipande).
 \emph{Journal of the American Mathematical Society} 20(2), 499--549,
 2007. --- \emph{The MNOP conjecture and its proof for toric
 threefolds.}
\item
 \textbf{``Seiberg--Witten prepotential from instanton counting''}
 (with N. Nekrasov). \emph{Advances in Theoretical and Mathematical
 Physics} 7(5), 831--864, 2003. --- \emph{Proof of Nekrasov's
 conjecture relating the instanton partition function to the
 Seiberg--Witten prepotential.}
\item
 \textbf{``Gromov--Witten theory, Hurwitz numbers, and matrix models''}
 (with R. Pandharipande). \emph{Annals of Mathematics} 153(1),
 189--228, 2001. --- \emph{The ELSV formula and its proof using
 Gromov--Witten theory.}
\item
 \textbf{``Quantum cohomology of the Hilbert scheme of points on a
 surface''} (with D. Maulik). \emph{Inventiones Mathematicae} 209(2),
 459--547, 2017. --- \emph{Computation of the quantum cohomology ring
 of Hilbert schemes using stable envelopes.}
\item
 \textbf{``Random tilings and the frozen boundary''} (with R. Kenyon,
 S. Sheffield). \emph{Journal of the American Mathematical Society}
 19(4), 899--951, 2006. --- \emph{The frozen boundary of random lozenge
 tilings and its description by algebraic curves.}
\item
 \textbf{``A new method for computing the volume of a line bundle.''}
 \emph{Duke Mathematical Journal} 101(3), 441--467, 2000. ---
 \emph{Introduction of Okounkov bodies and their application to volumes
 of line bundles.}
\item
 \textbf{``The Gromov--Witten/Donaldson--Thomas correspondence for
 toric threefolds''} (with R. Pandharipande). \emph{Inventiones
 Mathematicae} 178(2), 319--386, 2009. --- \emph{Complete proof of the
 GW/DT correspondence for toric varieties.}
\item
 \textbf{``Stable envelopes and quantum groups''} (with D. Maulik).
 \emph{Pure and Applied Mathematics Quarterly} 16(4), 1133--1205, 2020.
 --- \emph{Geometric construction of R-matrices using stable
 envelopes.}
\end{enumerate}

\begin{center}\rule{0.5\linewidth}{0.5pt}\end{center}

\section{19. Ongoing Work and Future
Directions}\label{ch2006andreiokounkov:ongoing-work-and-future-directions}

\subsection{19.1 Stable Envelopes and Geometric Representation
Theory}\label{ch2006andreiokounkov:stable-envelopes-and-geometric-representation-theory}

Okounkov continues to develop his theory of stable envelopes and its
applications to geometric representation theory. Current directions
include: - Elliptic stable envelopes and their relationship to elliptic
quantum groups. - Application of stable envelopes to the study of 3D
mirror symmetry. - Connections between stable envelopes and the theory
of quantum integrable systems.

\subsection{19.2 Moduli Spaces and Enumerative
Geometry}\label{ch2006andreiokounkov:moduli-spaces-and-enumerative-geometry}

He continues to work on enumerative geometry of moduli spaces,
including: - The Gromov--Witten theory of Calabi--Yau threefolds. - The
relationship between enumerative geometry and the representation theory
of vertex algebras. - Wall-crossing formulas in Donaldson--Thomas
theory.

\subsection{19.3 Random Structures and Integrable
Probability}\label{ch2006andreiokounkov:random-structures-and-integrable-probability}

The study of random surfaces, random partitions, and their limit shapes
continues. Current work focuses on: - The connection between random
tilings and the Gaussian free field. - The relationship between random
matrix theory and the representation theory of infinite-dimensional
groups.

\subsection{19.4 Connections to Physics}\label{ch2006andreiokounkov:connections-to-physics}

Okounkov remains interested in the connections between his work and
theoretical physics, including: - The AGT correspondence and its
generalizations. - Supersymmetric gauge theories and their moduli
spaces. - The relationship between topological string theory and
enumerative geometry.

\begin{center}\rule{0.5\linewidth}{0.5pt}\end{center}

\section{20. Conclusion}\label{ch2006andreiokounkov:conclusion}

Andrei Okounkov is one of the most versatile and influential
mathematicians of his generation. His work spans probability,
representation theory, algebraic geometry, and mathematical physics, and
he has made fundamental contributions to each field. The Fields Medal
citation --- ``for his contributions to bridging probability,
representation theory and algebraic geometry'' --- captures the essence
of his mathematics: he builds bridges between fields, importing ideas
from one to solve problems in another.

The Okounkov--Vershik approach transformed the representation theory of
symmetric groups. The MNOP conjecture and its proof transformed
enumerative geometry. The Nekrasov--Okounkov formula provided rigorous
foundations for instanton counting in supersymmetric gauge theory.
Okounkov bodies became a fundamental tool in birational geometry. The
frozen boundary of random tilings connected probability, algebraic
geometry, and statistical mechanics.

At Columbia University, Okounkov continues to push the boundaries of
mathematics. His work on stable envelopes and the geometric construction
of R-matrices is opening new directions in geometric representation
theory. His ability to see connections between distant fields --- and to
develop the technical machinery needed to exploit those connections ---
ensures that his influence will continue to grow.

\begin{center}\rule{0.5\linewidth}{0.5pt}\end{center}

\begin{center}\rule{0.5\linewidth}{0.5pt}\end{center}

\section*{Exercises}
\addcontentsline{toc}{section}{Exercises}

\begin{exercise} Compute the character table of the symmetric group $S_4$. \end{exercise}

\begin{exercise} Verify the Weyl dimension formula for the irreducible representation of $\\mathfrak{sl}_2$ with highest weight $n$. \end{exercise}

\begin{exercise} Show that the Plancherel measure on partitions of $n$ converges to the semicircle distribution as $n \\to \\infty$. \end{exercise}

\begin{exercise} Compute the correlation kernel for the discrete sine kernel and verify it reproduces the GUE distribution. \end{exercise}

\begin{center}\rule{0.5\linewidth}{0.5pt}\end{center}

\section*{Summary}
\addcontentsline{toc}{section}{Summary}

Okounkov established deep connections between representation theory, mathematical physics, and probability. His work on random matrices, the representation theory of symmetric groups, and the connection to the Riemann zeta function revealed unexpected links between algebra, analysis, and number theory.

\section{21. References and Further
Reading}\label{ch2006andreiokounkov:references-and-further-reading}

\subsection{Official Sources}\label{ch2006andreiokounkov:official-sources}

\begin{itemize}
\tightlist
\item
 IMU Fields Medal 2006 Announcement:
 https://www.mathunion.org/imu-awards/fields-medal/fields-medal-2006
\item
 Fields Medal Citation for Andrei Okounkov:
 https://www.mathunion.org/fileadmin/IMU/Prizes/Fields/2006/okounkov-citation.pdf
\item
 Columbia University: ``Andrei Okounkov awarded the Fields Medal.''
 2006.
\end{itemize}

\subsection{Biographical and Overview
Articles}\label{ch2006andreiokounkov:biographical-and-overview-articles}

\begin{itemize}
\tightlist
\item
 Borodin, A. ``The Work of Andrei Okounkov.'' ICM 2006 Laudatio.
\item
 Klarreich, E. ``Bridging Fields --- Andrei Okounkov's Mathematical
 Connections.'' \emph{Science News}, 2006.
\item
 Okounkov, A. ``Mathematics: The Science of Patterns.''
 \emph{Proceedings of the ICM}, 2006.
\end{itemize}

\subsection{Key Technical References}\label{ch2006andreiokounkov:key-technical-references}

\begin{itemize}
\tightlist
\item
 Okounkov, A. and Vershik, A. ``A new approach to the representation
 theory of the symmetric groups.'' \emph{Selecta Math.} 2 (1996),
 581--605.
\item
 Maulik, D., Nekrasov, N., Okounkov, A., and Pandharipande, R.
 ``Gromov--Witten theory and Donaldson--Thomas theory I.''
 \emph{Compositio Math.} 142 (2006), 1263--1285.
\item
 Maulik, D., Nekrasov, N., Okounkov, A., and Pandharipande, R.
 ``Gromov--Witten theory and Donaldson--Thomas theory II.'' \emph{JAMS}
 20 (2007), 499--549.
\item
 Nekrasov, N. and Okounkov, A. ``Seiberg--Witten prepotential from
 instanton counting.'' \emph{Adv. Theor. Math. Phys.} 7 (2003),
 831--864.
\item
 Kenyon, R., Okounkov, A., and Sheffield, S. ``Dimers and the frozen
 boundary.'' \emph{JAMS} 19 (2006), 899--951.
\item
 Okounkov, A. ``A new method for computing the volume of a line
 bundle.'' \emph{Duke Math. J.} 101 (2000), 441--467.
\item
 Okounkov, A. ``Random partitions and the frozen boundary.''
 \emph{Current Developments in Mathematics}, 2005.
\end{itemize}

\begin{center}\rule{0.5\linewidth}{0.5pt}\end{center}


\begin{center}\rule{0.5\linewidth}{0.5pt}\end{center}

\section*{See Also}
\addcontentsline{toc}{section}{See Also}

This chapter relates to the following work by other Fields Medalists:
\begin{itemize}
  \item \textbf{\hyperref[ch2022marynaviazovska:maryna-viazovska-fields-medal-2022-a-comprehensive-report-on-her-work-and-impact]{See Chapter on related work}} --- Viazovska's work on sphere packing relates.
  \item \textbf{\hyperref[ch2018caucherbirkar:caucher-birkar-fields-medal-2018-a-comprehensive-report-on-his-work-and-impact]{See Chapter on related work}} --- Birkar's work on birational geometry relates.
  \item \textbf{\hyperref[ch1998maximkontsevich:maxim-kontsevich-fields-medal-1998-a-comprehensive-report-on-his-work-and-impact]{See Chapter on related work}} --- Okounkov's work on moduli spaces relates.
\end{itemize}

\emph{(End of report)}
